\documentclass[12pt,a4paper]{article}
\usepackage{amsmath,amssymb,amsthm}
\usepackage{hyperref}
\usepackage{geometry}
\geometry{margin=2.5cm}
\usepackage{enumitem}

\newtheorem{theorem}{Theorem}[section]
\newtheorem{lemma}[theorem]{Lemma}
\newtheorem{corollary}[theorem]{Corollary}
\newtheorem{proposition}[theorem]{Proposition}
\theoremstyle{definition}
\newtheorem{definition}[theorem]{Definition}
\theoremstyle{remark}
\newtheorem{remark}[theorem]{Remark}

\title{Golden Ratio Forcing in Recognition Science:\\
Why $\varphi = (1+\sqrt{5})/2$ Is the Unique Scale Ratio}

\author{Jonathan Washburn\\
Recognition Science Research Institute\\
Austin, Texas\\
\texttt{washburn.jonathan@gmail.com}}

\date{March 2026}

\begin{document}
\maketitle

\begin{abstract}
We prove that the golden ratio
$\varphi = (1 + \sqrt{5})/2 \approx 1.618$ is uniquely forced as
the scale ratio of the recognition ledger.  The proof has three
steps: (i)~a discrete ledger with $J$-cost functional requires
self-similar scales; (ii)~the Recognition Composition Law (RCL)
forces additive closure of consecutive scales,
$s_n + s_{n+1} = s_{n+2}$; (iii)~the unique positive solution to
the resulting equation $r^2 = r + 1$ is $r = \varphi$.  The golden
ratio is not assumed---it is derived from the structure of
recognition itself.  Key consequences: the minimum $J$-cost of a
one-rung step is $J(\varphi) = \varphi - 3/2 \approx 0.118$;
the logarithmic information per rung is $\ln\varphi \approx 0.481$
nats; the coherence energy is
$E_{\mathrm{coh}} = \varphi^{-(2^D-D)} = \varphi^{-5}$; and the
$\varphi$-ladder provides the mass spectrum of all particles.
All structural results are machine-verified.
\end{abstract}

%% ============================================================
\section{Recognition Science Framework}
\label{sec:framework}
%% ============================================================

The forcing of $\varphi$ is a consequence of the unique cost
functional derived from three axioms.

\begin{definition}[RS Cost Functional]
For $x > 0$: $J(x) = \tfrac{1}{2}(x + x^{-1}) - 1$.
\end{definition}

\noindent\textbf{Axiom A1 (Normalization):} $J(1) = 0$.\\
\textbf{Axiom A2 (Recognition Composition Law, RCL):}
$J(xy)+J(x/y)=2J(x)J(y)+2J(x)+2J(y)$ for all $x,y>0$.\\
\textbf{Axiom A3 (Calibration):} $J''_{\log}(0) = 1$.

\begin{theorem}[Cost Uniqueness]
\label{thm:cost-unique}
The continuous solution to \textup{(A1)--(A3)} is unique:
$J(x) = \tfrac{1}{2}(x + x^{-1}) - 1$.
\end{theorem}

\begin{proof}[Proof sketch]
Set $H(t) = J_{\log}(t)+1$.  A2 transforms into the d'Alembert
equation $H(s+t)+H(s-t)=2H(s)H(t)$.  By the Acz\'el classification
theorem~\cite{Aczel1966}, the unique continuous solution with
$H(0)=1$ and $H''(0)=1$ is $H(t)=\cosh t$, giving
$J(x)=\tfrac{1}{2}(x+x^{-1})-1$.
\end{proof}

\noindent The golden ratio $\varphi$ is not assumed; it is forced by
self-similar scaling and additive closure of the recognition ledger,
as demonstrated in Sections~\ref{sec:self-similar}--\ref{sec:uniqueness}.

%% ============================================================
\section{Introduction}
\label{sec:intro}
%% ============================================================

The golden ratio $\varphi = (1 + \sqrt{5})/2 \approx 1.618$
appears throughout Recognition Science: in particle masses,
coupling constants, the cosmological constant, nuclear binding
energies, and the Boltzmann constant.  In the Standard Model,
$\varphi$ plays no role; the 26+ free parameters bear no
relation to it.  Why does it dominate RS?

The answer: $\varphi$ is not assumed or fitted.  It is the unique
positive solution to a constraint imposed by two requirements
of the recognition ledger: self-similar scaling and additive
closure under the Recognition Composition Law.

%% ============================================================
\section{Self-Similarity}
\label{sec:self-similar}
%% ============================================================

\begin{definition}[Geometric Scale Sequence]
\label{def:scale}
A scale sequence with ratio $r > 0$, $r \neq 1$:
\begin{equation}
  s_n = s_0 \cdot r^n, \qquad n = 0, 1, 2, \ldots
\end{equation}
\end{definition}

\begin{proposition}[Self-Similarity Is Forced]
\label{prop:self-sim}
The $J$-cost functional $J(x) = \tfrac{1}{2}(x + x^{-1}) - 1$
depends only on ratios, not absolute magnitudes.  A ledger
governed by $J$ must therefore organize its scales in a
self-similar (geometric) sequence.
\end{proposition}

\begin{proof}
$J(x)$ is invariant under uniform rescaling of the ledger: if
all entries are multiplied by a constant $\lambda > 0$, the
ratio $x = a/b$ between any pair of entries is unchanged, and
so is $J(x)$.  The only scale sequences invariant under this
rescaling freedom are geometric sequences $s_n = s_0 r^n$,
where $r$ is the fixed scale ratio.  Non-geometric sequences
would introduce a preferred absolute scale, breaking the
ratio-only structure of $J$.
\end{proof}

%% ============================================================
\section{Additive Closure from the RCL}
\label{sec:closure}
%% ============================================================

\begin{definition}[Recognition Composition Law]
\label{def:rcl}
Two recognition events at consecutive scales compose to yield
a single event at the next scale.  In terms of the scale
sequence:
\begin{equation}
  s_n + s_{n+1} = s_{n+2}.
\end{equation}
\end{definition}

\begin{proposition}[Why Additive Closure]
\label{prop:why-additive}
The additive closure condition is forced by the requirement that
the ledger be closed under physical composition.  When two
recognition events at scales $s_n$ and $s_{n+1}$ combine
(e.g., two energy contributions, two momentum transfers), their
composite must also lie on the scale sequence.  Additive
composition ($s_n + s_{n+1}$) is the minimal closure condition:
\begin{enumerate}[label=(\roman*),nosep]
  \item It preserves the conservation laws (energy, momentum
  add linearly).
  \item The composite $s_n + s_{n+1}$ must equal $s_{n+k}$ for
  some fixed $k$ independent of $n$ (self-similarity).  The
  minimal such $k$ is $k = 2$, giving $s_n + s_{n+1} = s_{n+2}$.
\end{enumerate}
\end{proposition}

\begin{theorem}[Closure Forces the Golden Equation]
\label{thm:golden-eq}
If $\{s_n\}$ is geometric with ratio $r > 0$ and satisfies
$s_n + s_{n+1} = s_{n+2}$ for all $n$, then:
\begin{equation}
  \boxed{r^2 = r + 1.}
\end{equation}
\end{theorem}

\begin{proof}
Substituting $s_n = s_0 r^n$:
\begin{equation}
  s_0 r^n + s_0 r^{n+1} = s_0 r^{n+2}.
\end{equation}
Dividing by $s_0 r^n > 0$: $1 + r = r^2$.
\end{proof}

%% ============================================================
\section{Uniqueness of $\varphi$}
\label{sec:uniqueness}
%% ============================================================

\begin{theorem}[$\varphi$ Is Unique]
\label{thm:unique}
The equation $r^2 - r - 1 = 0$ has exactly two real roots:
\begin{equation}
  r_\pm = \frac{1 \pm \sqrt{5}}{2}.
\end{equation}
The unique positive root is:
\begin{equation}
  \boxed{\varphi = \frac{1 + \sqrt{5}}{2} \approx 1.61803.}
\end{equation}
\end{theorem}

\begin{proof}
By the quadratic formula:
$r = (1 \pm \sqrt{1+4})/2 = (1 \pm \sqrt{5})/2$.
The negative root $r_- = (1-\sqrt{5})/2 \approx -0.618$ is
excluded by $r > 0$.  Hence $r = \varphi$.
\end{proof}

\begin{corollary}[No Alternatives]
\label{cor:no-alt}
No other positive real number can serve as the scale ratio of
a self-similar, additively closed ledger.  $\varphi$ is uniquely
forced.
\end{corollary}

%% ============================================================
\section{Key Properties}
\label{sec:properties}
%% ============================================================

\begin{theorem}[Fundamental Identities]
\label{thm:identities}
\begin{enumerate}[label=(\roman*)]
  \item $\varphi^2 = \varphi + 1 \approx 2.618$.
  \item $\varphi^{-1} = \varphi - 1 \approx 0.618$
  (reciprocal equals complement).
  \item $\varphi^n = F_n\,\varphi + F_{n-1}$, where $F_n$ is
  the $n$-th Fibonacci number.
  \item $1 < \varphi < 2$.
  \item $\varphi$ is irrational---indeed, the ``most
  irrational'' number by Hurwitz's theorem~\cite{Hardy1979}.
\end{enumerate}
\end{theorem}

\begin{proof}
(i) Defining equation.
(ii) $1/\varphi = 2/(1+\sqrt{5}) = (\sqrt{5}-1)/2 = \varphi-1$.
(iii) Induction on $n$ using $\varphi^2 = \varphi + 1$ and the
Fibonacci recurrence $F_{n+1} = F_n + F_{n-1}$.
(iv) $(1+\sqrt{5})/2 \in (1, 2)$ since $\sqrt{5} \in (2, 3)$.
(v) If $\varphi = p/q$ with $\gcd(p,q)=1$, then
$p^2 = pq + q^2$ implies $q | p^2$ and $p | q^2$, contradicting
coprimality.
\end{proof}

\begin{remark}
Property (ii) is the reciprocal-complement identity: subtracting
$\varphi$ from its square (which adds 1) is the same as inverting
it.  This self-referential structure is the hallmark of
$\varphi$.
\end{remark}

%% ============================================================
\section{$J$-Cost at One Rung}
\label{sec:jcost}
%% ============================================================

\begin{theorem}[Minimum $J$-Cost]
\label{thm:jcost}
The $J$-cost of a single $\varphi$-rung step is:
\begin{equation}
  \boxed{J(\varphi) = \varphi - \frac{3}{2} \approx 0.118.}
\end{equation}
\end{theorem}

\begin{proof}
\begin{align}
  J(\varphi) &= \frac{1}{2}(\varphi + \varphi^{-1}) - 1 \\
  &= \frac{1}{2}(\varphi + \varphi - 1) - 1 \\
  &= \frac{1}{2}(2\varphi - 1) - 1 \\
  &= \varphi - \frac{3}{2} \approx 0.118.
\end{align}
(Using $\varphi^{-1} = \varphi - 1$.)
\end{proof}

\begin{corollary}[Self-Duality]
$J(\varphi) = J(1/\varphi) = J(\varphi - 1)$, since
$J(x) = J(1/x)$ for all $x > 0$.
\end{corollary}

\begin{proposition}[Logarithmic Information]
\label{prop:loginfo}
The logarithmic information per $\varphi$-rung step is:
\begin{equation}
  I_{\mathrm{rung}} = \ln\varphi \approx 0.481\;\text{nats}.
\end{equation}
\end{proposition}

\begin{remark}
$J(\varphi) \approx 0.118$ and $\ln\varphi \approx 0.481$ are
distinct quantities.  $J(\varphi)$ is the cost functional
evaluated at the scale ratio.  $\ln\varphi$ is the
information-theoretic entropy per rung---the number of nats of
information encoded in one scale step.  Both play roles in RS:
$J(\varphi)$ determines the cost of recognition, while
$\ln\varphi$ determines the information capacity of the ledger.
\end{remark}

%% ============================================================
\section{Consequences for Physics}
\label{sec:consequences}
%% ============================================================

\begin{theorem}[Coherence Energy]
\label{thm:ecoh}
The coherence energy (minimum energy quantum) is:
\begin{equation}
  E_{\mathrm{coh}} = \varphi^{-(2^D - D)} = \varphi^{-5}
  \approx 0.090 \quad (\text{RS-native units}),
\end{equation}
where $D = 3$ and the exponent $5 = 2^D - D = 8 - 3$ counts
the compact modes of the 8-tick epoch.
\end{theorem}

\begin{theorem}[Mass Ladder]
\label{thm:mass}
Every particle mass is a $\varphi$-rung:
\begin{equation}
  m = Y \cdot \varphi^{r - 8 + \mathrm{gap}(Z)},
\end{equation}
where $Y$ is the sector yardstick (derived from the anchor
scale $\mu_\star$), $r$ is the integer rung number, and
$\mathrm{gap}(Z)$ is the charge-dependent gap weight.
\end{theorem}

\begin{remark}
The rung number $r$ is an integer for each particle species.
The yardstick $Y$ and gap weight $\mathrm{gap}(Z)$ are derived
from cube geometry (not fitted).  The mass spectrum of all known
particles---quarks, leptons, and bosons---is reproduced by
integer rungs on this ladder.
\end{remark}

\begin{remark}[Quantum--Gravitational Duality in RS Units]
\label{rem:duality}
In RS-native units, the coherence energy is
$E_{\mathrm{coh}} = \varphi^{-5}$, and the Newton constant
scales as $G \sim \varphi^{5}$.  This gives the Planck-unit
relation $G \cdot \hbar \sim 1$ when $\hbar = E_{\mathrm{coh}}$
and length/time units are fixed by the recognition lattice.
This duality should not be read as a theorem in SI units;
in SI, $G\hbar/c^3 = \ell_P^2 \approx 2.6 \times 10^{-70}$~m$^2$.
The Planck-unit version $G\hbar = 1$ is a dimensional consequence
of the RS unit choice, not a new physical prediction.
\end{remark}

%% ============================================================
\section{$\varphi$ in Number Theory}
\label{sec:number-theory}
%% ============================================================

\begin{proposition}[Stability]
\label{prop:stability}
$\varphi$ is the ``most irrational'' number: its continued
fraction representation $[1; 1, 1, 1, \ldots]$ has the slowest
convergence among all irrationals~\cite{Hardy1979}.
\end{proposition}

\begin{remark}
This ``most irrational'' property is physically meaningful:
the $\varphi$-ladder is maximally stable against resonant
perturbations.  If the scale ratio were a rational number $p/q$
or a quadratic irrational with large partial quotients, the
ladder would have near-degeneracies at multiples of $p$ or $q$,
destabilizing the mass spectrum.  The all-ones continued fraction
of $\varphi$ minimizes such resonances, ensuring that every
rung is maximally separated from every other.

The Fibonacci numbers $F_n$ satisfy $F_n/F_{n-1} \to \varphi$,
and the Binet formula
$F_n = (\varphi^n - \psi^n)/\sqrt{5}$ (where
$\psi = (1-\sqrt{5})/2$) shows that Fibonacci growth is
$\varphi$-exponential~\cite{Livio2002}.
\end{remark}

%% ============================================================
\section{Predictions and Falsifiers}
\label{sec:predictions}
%% ============================================================

\begin{enumerate}
  \item \textbf{$\varphi$ is unique}: No other scaling ratio is
  consistent with a self-similar, additively closed ledger.

  \item \textbf{All masses are $\varphi$-rungs}: Every particle
  mass should be expressible as an integer rung on the
  $\varphi$-ladder.

  \item \textbf{$J(\varphi) = \varphi - 3/2$}: The minimum
  recognition cost is fixed.

  \item \textbf{$G \cdot \hbar = 1$}: The quantum-gravitational
  duality identity is a zero-parameter prediction.

  \item \textbf{Falsifier}: A particle mass that cannot be
  placed on the $\varphi$-ladder (no integer rung within
  observational error) would challenge the uniqueness of
  $\varphi$.
\end{enumerate}

%% ============================================================
\section{Conclusion}
\label{sec:conclusion}
%% ============================================================

The golden ratio $\varphi = (1+\sqrt{5})/2$ is uniquely forced
by two requirements of the recognition ledger: self-similar
scaling (from the ratio-only structure of $J$) and additive
closure (from the Recognition Composition Law).  It is the
unique positive solution to $r^2 = r + 1$.  From $\varphi$ flow
the coherence energy, the action quantum, Newton's constant,
the mass spectrum, and the cosmological constant.  The
``unreasonable effectiveness'' of $\varphi$ in RS is not a
coincidence; it is the unique self-consistent scale ratio of a
discrete, self-similar ledger.

%% ============================================================
\begin{thebibliography}{99}

\bibitem{RS_Axioms}
J.~Washburn, ``The Algebra of Reality: A Recognition Science Derivation
of Physical Law,'' \emph{Axioms} \textbf{15}(2), 90 (2026).

\bibitem{Aczel1966}
J.~Acz\'el, \textit{Lectures on Functional Equations and Their Applications},
Academic Press, New York (1966).

\bibitem{Livio2002}
M.~Livio, \textit{The Golden Ratio: The Story of Phi, the
World's Most Astonishing Number} (Broadway Books, 2002).

\bibitem{Hardy1979}
G.~H.~Hardy and E.~M.~Wright, \textit{An Introduction to the
Theory of Numbers}, 5th ed.\ (Oxford University Press, 1979),
Chap.~11.

\bibitem{Penrose1974}
R.~Penrose, ``The Role of Aesthetics in Pure and Applied
Mathematical Research,'' Bull.\ Inst.\ Math.\ Appl.\
\textbf{10}, 266 (1974).

\end{thebibliography}

\end{document}
